\documentclass[11pt, letterpaper]{article}
% ============================================================
% preamble.tex — Shared LaTeX preamble for Learning to Autolens
% ============================================================
% Contains: packages, physics notation macros, formatting.
% Include in each module via % ============================================================
% preamble.tex — Shared LaTeX preamble for Learning to Autolens
% ============================================================
% Contains: packages, physics notation macros, formatting.
% Include in each module via % ============================================================
% preamble.tex — Shared LaTeX preamble for Learning to Autolens
% ============================================================
% Contains: packages, physics notation macros, formatting.
% Include in each module via \input{../preamble}
% ============================================================

% --- Packages ---
\usepackage[utf8]{inputenc}
\usepackage[T1]{fontenc}
\usepackage{amsmath, amssymb, amsthm}
\usepackage{physics}       % \vb, \grad, \div, \curl, \dd, etc.
\usepackage{mathtools}     % \coloneqq, \prescript
\usepackage{bm}            % \bm for bold math
\usepackage{graphicx}
\usepackage{float}
\usepackage{booktabs}      % Professional tables
\usepackage{hyperref}
\usepackage{cleveref}      % \cref for smart cross-refs
\usepackage{xcolor}
\usepackage[margin=1in]{geometry}
\usepackage{enumitem}      % Custom list formatting
\usepackage{fancyhdr}      % Headers and footers
\usepackage{tcolorbox}     % Colored boxes for key results
\usepackage{listings}      % Code listings

% --- Color scheme ---
\definecolor{codeblue}{RGB}{31, 119, 180}
\definecolor{codegray}{RGB}{128, 128, 128}
\definecolor{keyresult}{RGB}{39, 123, 69}
\definecolor{exercise}{RGB}{178, 102, 0}

% --- tcolorbox environments ---
\newtcolorbox{keyresult}[1][]{
  colback=keyresult!5, colframe=keyresult!70,
  fonttitle=\bfseries, title={Key Result}, #1
}

\newtcolorbox{pythonbox}[1][]{
  colback=codeblue!5, colframe=codeblue!50,
  fonttitle=\bfseries, title={PyAutoLens Implementation}, #1
}

\newtcolorbox{exercisebox}[1][]{
  colback=exercise!5, colframe=exercise!50,
  fonttitle=\bfseries, title={Exercise}, #1
}

% --- Code listing style ---
\lstset{
  language=Python,
  basicstyle=\ttfamily\small,
  keywordstyle=\color{codeblue}\bfseries,
  commentstyle=\color{codegray}\itshape,
  stringstyle=\color{keyresult},
  breaklines=true,
  frame=single,
  framerule=0.5pt,
  rulecolor=\color{codegray!50},
}

% --- Physics macros ---
% Vectors (bold upright)
\newcommand{\vtheta}{\bm{\theta}}       % Image-plane position
\newcommand{\vbeta}{\bm{\beta}}         % Source-plane position
\newcommand{\valpha}{\bm{\alpha}}       % Deflection angle
\newcommand{\vx}{\bm{x}}               % Generic position

% Lensing quantities
\newcommand{\thetaE}{\theta_{\rm E}}    % Einstein radius
\newcommand{\SigmaCr}{\Sigma_{\rm cr}}  % Critical surface density
\newcommand{\Dd}{D_{\rm d}}             % Angular diameter distance to lens
\newcommand{\Ds}{D_{\rm s}}             % Angular diameter distance to source
\newcommand{\Dds}{D_{\rm ds}}           % Angular diameter distance lens-source

% Statistical quantities
\newcommand{\chisq}{\chi^2}             % Chi-squared
\newcommand{\chisqred}{\chi^2_{\rm red}} % Reduced chi-squared
\newcommand{\Like}{\mathcal{L}}         % Likelihood
\newcommand{\Evid}{\mathcal{Z}}         % Evidence
\newcommand{\Post}{\mathcal{P}}         % Posterior

% Abbreviations for references
\newcommand{\CK}{C\&K}                  % Congdon & Keeton (2018)
\newcommand{\NB}{N\&B}                  % Narayan & Bartelmann (1997)
\newcommand{\Night}{Nightingale+18}     % Nightingale, Dye & Massey (2018)

% --- Headers ---
\pagestyle{fancy}
\fancyhf{}
\fancyhead[L]{\small\textit{Learning to Autolens}}
\fancyhead[R]{\small\thepage}
\renewcommand{\headrulewidth}{0.4pt}

% --- Theorem environments ---
\theoremstyle{definition}
\newtheorem{definition}{Definition}[section]
\newtheorem{example}{Example}[section]

% ============================================================

% --- Packages ---
\usepackage[utf8]{inputenc}
\usepackage[T1]{fontenc}
\usepackage{amsmath, amssymb, amsthm}
\usepackage{physics}       % \vb, \grad, \div, \curl, \dd, etc.
\usepackage{mathtools}     % \coloneqq, \prescript
\usepackage{bm}            % \bm for bold math
\usepackage{graphicx}
\usepackage{float}
\usepackage{booktabs}      % Professional tables
\usepackage{hyperref}
\usepackage{cleveref}      % \cref for smart cross-refs
\usepackage{xcolor}
\usepackage[margin=1in]{geometry}
\usepackage{enumitem}      % Custom list formatting
\usepackage{fancyhdr}      % Headers and footers
\usepackage{tcolorbox}     % Colored boxes for key results
\usepackage{listings}      % Code listings

% --- Color scheme ---
\definecolor{codeblue}{RGB}{31, 119, 180}
\definecolor{codegray}{RGB}{128, 128, 128}
\definecolor{keyresult}{RGB}{39, 123, 69}
\definecolor{exercise}{RGB}{178, 102, 0}

% --- tcolorbox environments ---
\newtcolorbox{keyresult}[1][]{
  colback=keyresult!5, colframe=keyresult!70,
  fonttitle=\bfseries, title={Key Result}, #1
}

\newtcolorbox{pythonbox}[1][]{
  colback=codeblue!5, colframe=codeblue!50,
  fonttitle=\bfseries, title={PyAutoLens Implementation}, #1
}

\newtcolorbox{exercisebox}[1][]{
  colback=exercise!5, colframe=exercise!50,
  fonttitle=\bfseries, title={Exercise}, #1
}

% --- Code listing style ---
\lstset{
  language=Python,
  basicstyle=\ttfamily\small,
  keywordstyle=\color{codeblue}\bfseries,
  commentstyle=\color{codegray}\itshape,
  stringstyle=\color{keyresult},
  breaklines=true,
  frame=single,
  framerule=0.5pt,
  rulecolor=\color{codegray!50},
}

% --- Physics macros ---
% Vectors (bold upright)
\newcommand{\vtheta}{\bm{\theta}}       % Image-plane position
\newcommand{\vbeta}{\bm{\beta}}         % Source-plane position
\newcommand{\valpha}{\bm{\alpha}}       % Deflection angle
\newcommand{\vx}{\bm{x}}               % Generic position

% Lensing quantities
\newcommand{\thetaE}{\theta_{\rm E}}    % Einstein radius
\newcommand{\SigmaCr}{\Sigma_{\rm cr}}  % Critical surface density
\newcommand{\Dd}{D_{\rm d}}             % Angular diameter distance to lens
\newcommand{\Ds}{D_{\rm s}}             % Angular diameter distance to source
\newcommand{\Dds}{D_{\rm ds}}           % Angular diameter distance lens-source

% Statistical quantities
\newcommand{\chisq}{\chi^2}             % Chi-squared
\newcommand{\chisqred}{\chi^2_{\rm red}} % Reduced chi-squared
\newcommand{\Like}{\mathcal{L}}         % Likelihood
\newcommand{\Evid}{\mathcal{Z}}         % Evidence
\newcommand{\Post}{\mathcal{P}}         % Posterior

% Abbreviations for references
\newcommand{\CK}{C\&K}                  % Congdon & Keeton (2018)
\newcommand{\NB}{N\&B}                  % Narayan & Bartelmann (1997)
\newcommand{\Night}{Nightingale+18}     % Nightingale, Dye & Massey (2018)

% --- Headers ---
\pagestyle{fancy}
\fancyhf{}
\fancyhead[L]{\small\textit{Learning to Autolens}}
\fancyhead[R]{\small\thepage}
\renewcommand{\headrulewidth}{0.4pt}

% --- Theorem environments ---
\theoremstyle{definition}
\newtheorem{definition}{Definition}[section]
\newtheorem{example}{Example}[section]

% ============================================================

% --- Packages ---
\usepackage[utf8]{inputenc}
\usepackage[T1]{fontenc}
\usepackage{amsmath, amssymb, amsthm}
\usepackage{physics}       % \vb, \grad, \div, \curl, \dd, etc.
\usepackage{mathtools}     % \coloneqq, \prescript
\usepackage{bm}            % \bm for bold math
\usepackage{graphicx}
\usepackage{float}
\usepackage{booktabs}      % Professional tables
\usepackage{hyperref}
\usepackage{cleveref}      % \cref for smart cross-refs
\usepackage{xcolor}
\usepackage[margin=1in]{geometry}
\usepackage{enumitem}      % Custom list formatting
\usepackage{fancyhdr}      % Headers and footers
\usepackage{tcolorbox}     % Colored boxes for key results
\usepackage{listings}      % Code listings

% --- Color scheme ---
\definecolor{codeblue}{RGB}{31, 119, 180}
\definecolor{codegray}{RGB}{128, 128, 128}
\definecolor{keyresult}{RGB}{39, 123, 69}
\definecolor{exercise}{RGB}{178, 102, 0}

% --- tcolorbox environments ---
\newtcolorbox{keyresult}[1][]{
  colback=keyresult!5, colframe=keyresult!70,
  fonttitle=\bfseries, title={Key Result}, #1
}

\newtcolorbox{pythonbox}[1][]{
  colback=codeblue!5, colframe=codeblue!50,
  fonttitle=\bfseries, title={PyAutoLens Implementation}, #1
}

\newtcolorbox{exercisebox}[1][]{
  colback=exercise!5, colframe=exercise!50,
  fonttitle=\bfseries, title={Exercise}, #1
}

% --- Code listing style ---
\lstset{
  language=Python,
  basicstyle=\ttfamily\small,
  keywordstyle=\color{codeblue}\bfseries,
  commentstyle=\color{codegray}\itshape,
  stringstyle=\color{keyresult},
  breaklines=true,
  frame=single,
  framerule=0.5pt,
  rulecolor=\color{codegray!50},
}

% --- Physics macros ---
% Vectors (bold upright)
\newcommand{\vtheta}{\bm{\theta}}       % Image-plane position
\newcommand{\vbeta}{\bm{\beta}}         % Source-plane position
\newcommand{\valpha}{\bm{\alpha}}       % Deflection angle
\newcommand{\vx}{\bm{x}}               % Generic position

% Lensing quantities
\newcommand{\thetaE}{\theta_{\rm E}}    % Einstein radius
\newcommand{\SigmaCr}{\Sigma_{\rm cr}}  % Critical surface density
\newcommand{\Dd}{D_{\rm d}}             % Angular diameter distance to lens
\newcommand{\Ds}{D_{\rm s}}             % Angular diameter distance to source
\newcommand{\Dds}{D_{\rm ds}}           % Angular diameter distance lens-source

% Statistical quantities
\newcommand{\chisq}{\chi^2}             % Chi-squared
\newcommand{\chisqred}{\chi^2_{\rm red}} % Reduced chi-squared
\newcommand{\Like}{\mathcal{L}}         % Likelihood
\newcommand{\Evid}{\mathcal{Z}}         % Evidence
\newcommand{\Post}{\mathcal{P}}         % Posterior

% Abbreviations for references
\newcommand{\CK}{C\&K}                  % Congdon & Keeton (2018)
\newcommand{\NB}{N\&B}                  % Narayan & Bartelmann (1997)
\newcommand{\Night}{Nightingale+18}     % Nightingale, Dye & Massey (2018)

% --- Headers ---
\pagestyle{fancy}
\fancyhf{}
\fancyhead[L]{\small\textit{Learning to Autolens}}
\fancyhead[R]{\small\thepage}
\renewcommand{\headrulewidth}{0.4pt}

% --- Theorem environments ---
\theoremstyle{definition}
\newtheorem{definition}{Definition}[section]
\newtheorem{example}{Example}[section]


\fancyhead[C]{\small\textbf{Module 09: MGE \& Linear Light Profiles}}

\title{\textbf{Module 09: Multi-Gaussian Expansion \& Linear Light Profiles}\\[0.5em]
\large Theory Companion --- Learning to Autolens}
\author{Rodrigo C\'ordova Rosado\\
\small Harvard-Smithsonian Center for Astrophysics\\
\small \texttt{rodrigo.cordova\_rosado@cfa.harvard.edu}}
\date{\today}

\begin{document}
\maketitle

\begin{abstract}
This document develops the mathematical framework for linear light
profile fitting and Multi-Gaussian Expansions (MGEs) as used in
PyAutoLens.  We show that solving for light profile intensities is
a linear algebra problem identical in structure to the pixelized
source inversion of Module~05, derive the maximum-likelihood and
NNLS solutions, and explain why Gaussians form a particularly
efficient basis.
References: Emsellem, Monnet \& Bacon (1994); Cappellari (2002);
Lawson \& Hanson (1974).
\end{abstract}

\tableofcontents
\newpage

% ============================================================
\section{Linear Algebra of Light Profile Fitting}
\label{sec:linear-algebra}
% ============================================================

In a traditional lens model every light profile has an \texttt{intensity}
parameter that the non-linear search must explore.  A \emph{linear light
profile} removes intensity from the search: given fixed shape parameters,
the optimal intensity is found by solving a linear system.

\subsection{The Forward Model}

Suppose we have $K$ light profile components (e.g.\ bulge, disk, lensed
arcs), each with a known 2D shape but unknown intensity $a_k$.  Let
$\mathbf{b}_k$ be the PSF-convolved image of component $k$ evaluated at
unit intensity --- a column vector of $N_d$ image pixels.  Stack these
into the \textbf{blurring matrix}
$\mathbf{B} = [\mathbf{b}_1 \mid \cdots \mid \mathbf{b}_K]$, which has
dimensions $N_d \times K$.  The model image is then:

\begin{keyresult}[title={Forward Model}]
\begin{equation}
  \mathbf{d} = \mathbf{B}\,\mathbf{a} + \mathbf{n},
  \label{eq:forward}
\end{equation}
where $\mathbf{a} = (a_1,\ldots,a_K)^T$ is the intensity vector and
$\mathbf{n}$ is the noise with covariance
$\mathbf{C} = \mathrm{diag}(\sigma_1^2,\ldots,\sigma_{N_d}^2)$.
\end{keyresult}

This is \emph{exactly} the same structure as the pixelized source
inversion from Module~05 (\cref{eq:forward} $\leftrightarrow$
Module~05 eq.~1), with the replacement
\[
  \underbrace{\mathbf{M}}_{\text{lensing+PSF}} \mathbf{s}
  \;\longleftrightarrow\;
  \underbrace{\mathbf{B}}_{\text{PSF-convolved basis}} \mathbf{a}.
\]
The source pixels $\mathbf{s}$ are replaced by intensities $\mathbf{a}$,
and the lensing operator by the basis image matrix.  All the linear
algebra carries over.

\subsection{Maximum-Likelihood Solution}

Minimizing $\chisq(\mathbf{a}) = (\mathbf{d} - \mathbf{B}\mathbf{a})^T
\mathbf{C}^{-1}(\mathbf{d} - \mathbf{B}\mathbf{a})$ and setting
$\partial\chisq/\partial\mathbf{a} = 0$ gives the weighted
least-squares solution:

\begin{keyresult}[title={ML Intensity Solution}]
\begin{equation}
  \hat{\mathbf{a}}_{\rm ML}
  = \bigl(\mathbf{B}^T \mathbf{C}^{-1} \mathbf{B}\bigr)^{-1}
    \mathbf{B}^T \mathbf{C}^{-1} \mathbf{d}
  \label{eq:ml-solution}
\end{equation}
\end{keyresult}

\noindent The matrix $\mathbf{B}^T\mathbf{C}^{-1}\mathbf{B}$ is only
$K \times K$ (typically $K \sim 20$--60 for an MGE), so the solve is
trivially cheap.  Unlike the pixelized source case, \emph{no
regularization is needed} because $K \ll N_d$ and the system is
well-conditioned.

% ============================================================
\section{Positive-Only Constraint (NNLS)}
\label{sec:nnls}
% ============================================================

The unconstrained solution~\eqref{eq:ml-solution} can yield negative
$a_k$.  While mathematically valid, a negative surface brightness is
unphysical for a stellar light component.

PyAutoLens therefore solves the \textbf{Non-Negative Least Squares}
(NNLS) problem by default:
\begin{equation}
  \hat{\mathbf{a}}_{\rm NNLS}
  = \arg\min_{\mathbf{a} \ge 0}
    \bigl\|\mathbf{C}^{-1/2}(\mathbf{d} - \mathbf{B}\mathbf{a})\bigr\|^2
  \label{eq:nnls}
\end{equation}

The standard algorithm is that of Lawson \& Hanson (1974): an active-set
method that iteratively partitions the indices into a ``free'' set
(non-zero coefficients) and a ``fixed'' set (zeroed coefficients),
solving the unconstrained sub-problem on the free set at each step.
Convergence is guaranteed in a finite number of iterations.

\begin{exercisebox}
Construct a toy 1D example with three Gaussian basis functions where
the unconstrained least-squares solution gives one negative amplitude.
Verify that NNLS sets it to zero and redistributes flux to the
remaining components.
\end{exercisebox}

% ============================================================
\section{Multi-Gaussian Expansion Theory}
\label{sec:mge-theory}
% ============================================================

\subsection{The MGE Ansatz}

A Multi-Gaussian Expansion represents a surface brightness profile as a
sum of 2D Gaussians:
\begin{equation}
  I(R) = \sum_{j=1}^{K} a_j \,
  \exp\!\left(-\frac{R^2}{2\sigma_j^2}\right)
  \label{eq:mge-ansatz}
\end{equation}
where $R$ is the elliptical radius appropriate to each component and
$\{\sigma_j\}$ are the Gaussian widths (dispersions).

Historically, MGEs were developed for stellar dynamical modeling
(Emsellem, Monnet \& Bacon 1994) and efficient fitting algorithms were
given by Cappellari (2002).  Their adoption in lens modeling is more
recent.

\subsection{Why Gaussians Are Special}

Three properties make Gaussians a uniquely convenient basis:

\begin{enumerate}[nosep]
  \item \textbf{Analytic PSF convolution.}  If the PSF is itself
    approximated by a Gaussian with dispersion $\sigma_p$, the
    convolution of a basis Gaussian ($\sigma_j$) with the PSF is another
    Gaussian:

    \begin{keyresult}[title={Gaussian Convolution Property}]
    \begin{equation}
      \sigma_{\rm eff} = \sqrt{\sigma_j^2 + \sigma_p^2}
      \label{eq:conv-property}
    \end{equation}
    The convolved image is a Gaussian with the same centre and
    ellipticity, but broadened dispersion $\sigma_{\rm eff}$.
    \end{keyresult}

    In practice, PyAutoLens convolves numerically with an arbitrary PSF
    kernel (not necessarily Gaussian), but the near-Gaussian shape of
    each basis image means the convolution is well-behaved.

  \item \textbf{Analytic deprojection.}  A 2D Gaussian deprojects to a
    3D Gaussian (with the same $\sigma$ along the projected axis),
    enabling analytic computation of the gravitational potential and
    forces --- critical for dynamical modeling (Emsellem+94 Sec.~3).

  \item \textbf{Computational speed.}  Evaluating $\exp(-R^2/2\sigma^2)$
    is fast and numerically stable; no special functions (incomplete
    gamma, hypergeometric) are needed, unlike S\'ersic profiles.
\end{enumerate}

% ============================================================
\section{MGE as a Basis Expansion}
\label{sec:mge-basis}
% ============================================================

\subsection{The Sigma Grid}

PyAutoLens constructs the Gaussian widths on a logarithmic grid spanning
from sub-pixel scales to the mask radius:
\[
  \sigma_j \in \bigl\{\sigma_{\min},\;
  \sigma_{\min}\cdot r,\;
  \sigma_{\min}\cdot r^2,\;
  \ldots,\;
  \sigma_{\max}\bigr\}
\]
where $r$ is the common ratio.  Typical values: $\sigma_{\min} \approx
0.01''$ (sub-pixel), $\sigma_{\max} \approx$ mask radius, and
$K \sim 20$--30 Gaussians per group.  The log-spacing ensures roughly
equal sensitivity per octave of angular scale.

\subsection{Groups and \texttt{gaussian\_per\_basis}}

For multi-component morphologies (e.g.\ bulge + disk, or lens galaxy +
lensed arcs), the Gaussians are divided into \textbf{groups}.  Each
group shares:
\begin{itemize}[nosep]
  \item a common centre (\texttt{centre}),
  \item common ellipticity components (\texttt{ell\_comps}).
\end{itemize}

Within each group, the $K_g$ Gaussian widths $\{\sigma_j\}$ are fixed
(non-linear parameters would be needed only to adjust the grid, which
is not done), and the $K_g$ intensities $\{a_j\}$ are solved linearly.
The only non-linear parameters per group are:
\[
  \underbrace{x_c,\; y_c}_{\text{centre}} \;+\;
  \underbrace{e_1,\; e_2}_{\text{ellipticity}}
  \;=\; 4 \text{ parameters/group.}
\]

\begin{pythonbox}
A two-group MGE (lens light + source light) is set up as:
\begin{lstlisting}
# Group 1: lens galaxy (30 Gaussians)
lens_bulge = af.Model(
    al.lp_linear.Gaussian,
    sigma=af.LogUniformPrior(
        lower_limit=0.01, upper_limit=10.0
    ),
)
# Group 2: source galaxy (20 Gaussians)
source = af.Model(
    al.lp_linear.Gaussian,
    sigma=af.LogUniformPrior(
        lower_limit=0.01, upper_limit=5.0
    ),
)
\end{lstlisting}
In practice, the \texttt{al.util.model.mge\_model\_from()} utility
constructs the full model automatically (see Sec.~\ref{sec:reduction}).
\end{pythonbox}

% ============================================================
\section{Parameter Space Reduction}
\label{sec:reduction}
% ============================================================

The central advantage of MGE + linear solving is the dramatic reduction
in the dimensionality of the non-linear search.

\begin{table}[h]
\centering
\caption{Non-linear parameter count comparison for a bulge + disk
lens light model.}
\label{tab:param-count}
\begin{tabular}{lcc}
\toprule
\textbf{Approach} & \textbf{Non-linear params} & \textbf{Linear params} \\
\midrule
S\'ersic bulge + S\'ersic disk & 14 & 0 \\
\quad (each: $x, y, e_1, e_2, R_e, n, I$) & & \\
\addlinespace
MGE 2-group (30+30 Gaussians) & $2 \times 4 = 8$ & 60 \\
\quad (each group: $x, y, e_1, e_2$) & & \\
\addlinespace
MGE 1-group (shared geometry) & 4 & 30 \\
\bottomrule
\end{tabular}
\end{table}

\noindent The trade-off:
\begin{itemize}[nosep]
  \item \textbf{Fewer non-linear dimensions} $\Rightarrow$ the sampler
    converges faster and explores the posterior more efficiently.
  \item \textbf{Each likelihood evaluation is more expensive}: computing
    $K$ basis images and solving the $K \times K$ linear system adds
    overhead compared to evaluating a single S\'ersic profile.
\end{itemize}
In practice the dimensionality reduction dominates, and MGE models
typically complete 5--10$\times$ faster than equivalent S\'ersic models.

\begin{pythonbox}[title={The \texttt{mge\_model\_from} Utility}]
\begin{lstlisting}
import autolens as al
from autolens.util.model import mge_model_from

# Build an MGE model for the lens light
lens_bulge = mge_model_from(
    model=al.lp_linear.Gaussian,
    gaussian_per_basis=30,
)
# Returns an af.Collection of 30 Gaussian
# models sharing centre and ell_comps,
# with log-spaced sigma values.
\end{lstlisting}
\end{pythonbox}

% ============================================================
\section{MGE in the SLaM Pipeline}
\label{sec:slam}
% ============================================================

The standard SLaM (Source, Light, Mass) pipeline has four stages when
using pixelized sources (Module~04).  When MGE linear profiles are used
for \emph{both} lens light and source light, the pipeline simplifies to
three stages:

\begin{enumerate}[nosep]
  \item \textbf{SOURCE LP} --- Fit a parametric (MGE) source with a
    simple mass model (SIE + shear).  The source MGE intensities are
    solved linearly; the non-linear search explores only mass parameters
    and source geometry (centre, ellipticity).
  \item \textbf{LIGHT LP} --- Fix the mass model from Stage~1 and fit
    the lens galaxy light with an MGE.  This decouples the lens light
    from the mass, allowing a flexible multi-Gaussian representation
    without destabilizing the mass model.
  \item \textbf{MASS TOTAL} --- Fit the total mass model (e.g.\
    power-law ellipsoid + shear) with lens light and source light
    initialised from previous stages.
\end{enumerate}

\noindent \textbf{Why no SOURCE PIX stage?}  The pixelized source stage
exists to capture irregular source morphology that parametric profiles
miss.  When the source is already represented by an MGE with $\sim$20--30
Gaussians, it has enough flexibility to represent most source
morphologies without a pixelized reconstruction.  The SOURCE~PIX stage
is needed only when the source has structure on scales finer than the
densest Gaussian in the grid (e.g.\ highly irregular merging galaxies).

\subsection{Decision Guide}

\begin{table}[h]
\centering
\caption{When to use each source representation.}
\label{tab:decision}
\begin{tabular}{lp{8cm}}
\toprule
\textbf{Method} & \textbf{Best for} \\
\midrule
S\'ersic & Quick initial models; sources with known smooth morphology;
  when only 1--2 light components are needed. \\
\addlinespace
MGE (linear) & Default choice for most lenses; multi-component
  morphologies; fast convergence when the non-linear parameter
  budget is tight. \\
\addlinespace
Pixelized & Highly irregular or merging sources; when residuals from
  parametric fits show systematic structure; substructure studies. \\
\bottomrule
\end{tabular}
\end{table}

\begin{exercisebox}
Take the SLaM pipeline from Module~04 and replace the SOURCE~PIX stage
with a SOURCE~LP stage using an MGE source.  Compare the Bayesian
evidence of the two approaches on the same dataset.
\end{exercisebox}

% ============================================================
\section*{References}
% ============================================================

\begin{itemize}[nosep, leftmargin=*]
  \item Cappellari, M.\ (2002), MNRAS, 333, 400
  \item Emsellem, E., Monnet, G.\ \& Bacon, R.\ (1994), A\&A, 285, 723
  \item Lawson, C.L.\ \& Hanson, R.J.\ (1974),
    \textit{Solving Least Squares Problems}, Prentice-Hall
  \item Nightingale, J.W., Dye, S.\ \& Massey, R.J.\ (2018), MNRAS, 478, 4738
\end{itemize}

\end{document}
